% ============================================================
% QUESTION TYPE: TRUE / FALSE
%
% PEANUT EXAM TRUE/FALSE LAYOUT STANDARD
%
% AI INSTRUCTION:
% - Follow this format EXACTLY for True/False questions
% - Use \subsection*{Problem X} to start the section
% - Statements are listed using enumerate
% - Each statement begins with a fixed-width underline for the answer
% - Answers must be written as TRUE or FALSE
% - Control spacing using \vspace{...} only
% - Do NOT convert to MCQ or add explanations
% ============================================================

\newpage
\subsection*{Problem X}

\noindent For each statement below, write \textbf{TRUE} or \textbf{FALSE}
in the space provided.

\vspace{0.5cm}

\begin{enumerate}[label=\textbf{X.\arabic*}, left=0pt]

\item \underline{\hspace{3cm}}  
This is a factual or conceptual statement related to the course material.
It should be unambiguous and test conceptual understanding.

\vspace{0.3cm}

\item \underline{\hspace{3cm}}  
This statement may include inline mathematics such as $f(x)$ or $x_t$,
or reference standard definitions or properties.

\vspace{0.3cm}

\item \underline{\hspace{3cm}}  
The wording should be precise. Avoid opinion-based or vague statements.
Negations such as ``not'' may be used carefully.

\vspace{0.3cm}

\item \underline{\hspace{3cm}}  
Statements may refer to algorithms, models, assumptions, or limitations
discussed in the course.

\vspace{0.3cm}

\item \underline{\hspace{3cm}}  
A statement is either entirely correct or entirely incorrect
given standard textbook assumptions.

\vspace{0.3cm}

\end{enumerate}
